\section{Introduction}\label{sec:intro}
The Tunisian utility, the Société Tunisienne de l'Électricité et du Gaz (STEG), is deploying smart
meters. Press reports of a STEG briefing put the count at about 80,000 low-voltage (LV) meters in
pilot zones by the end of March 2026, with a first phase of 450,000 electricity meters
\cite{africanmanager2026compteurs,webdo2026compteurs}. STEG serves about 4.7 million subscribers,
most of them residential LV customers \cite{lapresse2025steg}. In the summer of 2026 the same network
went through rotating power cuts during a heatwave
\cite{newarab2026heatwave,africanmanager2026heat,webdo2026delestages}. These facts come from
secondary sources; they motivate the problem and enter no result.

As this advanced metering infrastructure (AMI) grows, operators will see groups of meters go silent.
From the meter data alone they cannot tell an \emph{outage}, which needs a crew, from a
\emph{communication failure}, which does not. A wrong call either wastes a truck roll or leaves
customers without power and unattended. SHAHID addresses this one question for one LV neighbourhood
under one medium-voltage/low-voltage (MV/LV) transformer, with 15-min readings. It is a probabilistic
artificial intelligence (AI) trust layer between smart-meter data and dispatch decisions.

\textbf{Hypothesis.} A meter on the LV side of the transformer still measures the energy of customers
whose data stopped, and no longer measures the energy of customers who lost power. If this difference
is large compared with the uncertainty of the expected consumption, the two causes can be separated.

\textbf{Contributions.} All results are SIMULATED.
\begin{enumerate}[leftmargin=1.4em]
\item A transformer-side energy-balance test with abstention that, with thresholds frozen on one
network, called no labelled communication failure an outage in every registered evaluation, on four
networks never used for tuning. On data with no events, random single-meter data losses still produced
about 15 to 33 false outage verdicts per transformer-year under the registered configuration; a post-hoc
gate reduced this to 0.2 to 1.0 (Section~\ref{sec:results}).
\item A measured operating envelope: the resolution limit, the reporting interval, and five conditions
under which the method weakens or fails (Section~\ref{sec:weakens}).
\item Governance measured, not assumed: an escalation table whose automatic rule was found to close
real outages and was made provisional, operator load per transformer-day, a hash-chained decision log
and an operator console preview (Sections~\ref{sec:explain}--\ref{sec:governance}).
\item A failure mechanism measured and removed: a powered meter merged into an outage group can
explain the whole residual; splitting silent groups by feeder removed the confident wrong verdicts it
caused, 4 to 0 (Section~\ref{sec:phase2}).
\item A verifiable evidence chain: protocols committed before each confirmatory run, negative results
reported, and every number in this paper checked by code against a claims ledger
(Section~\ref{sec:protocol}).
\end{enumerate}
